\documentclass[12pt]{article}

\usepackage{sbc-template}
\usepackage{graphicx,url}
\usepackage[utf8]{inputenc}
\usepackage[brazil]{babel}
\usepackage{amsmath}
\usepackage{booktabs}   % tabelas limpas (\toprule, \midrule, \bottomrule)
\usepackage{multirow}   % celulas mescladas em tabelas
\usepackage{float}      % posicionamento de figuras/tabelas
\usepackage{placeins}   % \FloatBarrier para conter floats por secao
\usepackage{hyperref}   % links clicaveis

\sloppy

\title{Classificação Multiclasse de Lesões da Mucosa Oral com Aprendizado Profundo em Imagens de Smartphone}

\author{
  Gustavo Novack Viana Lima,
  Henrick de Souza Silva,\\
  André Moura Koraleski,
  Caio Lucca dos Santos Oliveira
}

\address{
  Instituto de Informática -- Universidade Federal de Goiás (UFG)\\
  Goiânia -- GO -- Brasil
  \email{\{novack, henrick.henrick, andre.andre2, caiolucca\}@discente.ufg.br}
}

\begin{document}

\maketitle

\begin{abstract}
  Oral cancer diagnosis remains late-stage in most cases, limiting
  five-year survival to below 50\%.
  This work benchmarks three YOLO classification architectures
  (YOLOv8m-cls, YOLOv11m-cls, YOLO26m-cls) on the SMART-OM dataset —
  2,469 annotated intraoral smartphone images across four diagnostic
  classes — against published deep learning baselines.
  Inverse-frequency loss weighting and oversampling address severe class
  imbalance (86.9\% Normal).
  YOLO26m-cls surpasses the best baseline in macro sensitivity on the
  four-class task (69.57\% vs.\ 61.57\%, ResNet50) and on binary
  anomaly detection (66.15\% vs.\ 65.62\%, ResNet18).
\end{abstract}

\begin{resumo}
  O diagnóstico do câncer oral ocorre predominantemente em estágio
  avançado, com sobrevida em cinco anos inferior a 50\%.
  Este trabalho avalia três arquiteturas YOLO de classificação
  (YOLOv8m-cls, YOLOv11m-cls, YOLO26m-cls) no dataset SMART-OM —
  2.469 imagens intraorais de smartphone em quatro classes diagnósticas
  — frente aos baselines publicados, com ponderação da loss por
  frequência inversa e \textit{oversampling} para tratar o severo
  desbalanceamento (86,9\% Normal).
  O YOLO26m-cls supera o melhor baseline em sensibilidade macro na
  tarefa multiclasse (69,57\% vs.\ 61,57\%, ResNet50) e na detecção
  binária de anomalias (66,15\% vs.\ 65,62\%, ResNet18).
\end{resumo}

% =====================================================================
\section{Introdução}
\label{sec:introducao}

O câncer oral (CO) representa um grave problema de saúde pública global.
Segundo dados do GLOBOCAN 2020, cânceres de lábio e cavidade oral foram
responsáveis por 177.757 mortes e 377.713 novos casos em todo o mundo,
figurando como o 16° câncer mais incidente globalmente~\cite{globocan2020}.
A sobrevida em cinco anos permanece inferior a 50\%, pois a maioria dos
diagnósticos ocorre em estágios avançados da doença.
Condições precursoras — Desordens Potencialmente Malignas Orais (OPMD,
do inglês \textit{Oral Potentially Malignant Disorders}) — representam
uma janela terapêutica crítica: a detecção precoce e a diferenciação de
tecido saudável possibilita intervenção antes da progressão maligna.

O diagnóstico convencional baseia-se na inspeção visual por especialistas
em patologia oral, processo inerentemente subjetivo e suscetível a
variabilidade intra e interobservador.
Em países de baixa e média renda, a concentração de especialistas nos
grandes centros urbanos limita o acesso à triagem, resultando em
diagnósticos tardios~\cite{globocan2020}.
A maior parte dos \textit{frameworks} de visão computacional para lesões
orais foi desenvolvida sobre imagens histopatológicas — que exigem biópsia
e microscopia, incompatíveis com a atenção primária.
Imagens capturadas por smartphones surgem como alternativa escalável e de
baixo custo, mas a limitada disponibilidade de bases de dados públicas
nessa modalidade restringiu o avanço da área.

Nesse contexto, Kumar et al.~\cite{smartom} introduziram o SMART-OM
(\textit{SMARTphone-based dataset of Oral Mucosa images}), o único dataset
público de imagens intraorais identificado entre 332 disponíveis na área:
2.469 fotografias de 331 pacientes em quatro classes diagnósticas (Normal,
Variação do Normal, OPMD e Câncer Oral), coletadas com smartphones Android
e iOS em condições clínicas reais.
Os autores estabeleceram um benchmark avaliando seis arquiteturas —
ResNet18, ResNet34, ResNet50, VGG16, EfficientNet-b0 e Vision Transformer
(ViT)~\cite{resnet,efficientnet,vit} — em detecção binária de anomalias e
classificação multiclasse, atingindo 65,62\% de sensibilidade para a classe
Anormal (tarefa binária) e 61,57\% de sensibilidade macro (multiclasse).

O presente trabalho propõe e avalia três arquiteturas YOLO no modo de
classificação sobre o SMART-OM — YOLOv8m-cls, YOLOv11m-cls e
YOLO26m-cls~\cite{yolov8,yolo11,yolo26} — como alternativa às CNNs
clássicas e ao ViT do benchmark original.
Como segunda contribuição, propomos um pipeline em dois estágios:
detector binário de anomalias (Estágio~1) seguido de classificador de
três classes (Estágio~2), com ponderação da função de perda por frequência
inversa para mitigar o severo desbalanceamento do dataset.

\section{Formulação do Problema}
\label{sec:formulacao}

\subsection{Definição da Tarefa}

Dado um conjunto de imagens intraorais, o objetivo é aprender uma função
$f_\theta$ que mapeie cada fotografia da mucosa oral a uma categoria
diagnóstica. Avaliamos duas formulações complementares desta tarefa:

\textbf{T\textsubscript{1} — Detecção binária de anomalias.}
Classificação em duas classes — \textit{Normal} e \textit{Anormal} —
onde a segunda agrupa todas as imagens com alterações da mucosa
(Variação do Normal, OPMD e Câncer Oral).
O objetivo primário é maximizar a sensibilidade para a classe Anormal,
minimizando falsos negativos: clinicamente, deixar escapar uma lesão
tem consequências mais graves do que um falso alarme.

\textbf{T\textsubscript{2} — Classificação multiclasse.}
Classificação em quatro classes: \textit{Normal}, \textit{Variação do
Normal}, \textit{OPMD} e \textit{Câncer Oral}, replicando diretamente
a tarefa do benchmark SMART-OM~\cite{smartom}.
Esta formulação fornece maior granularidade diagnóstica, permitindo
distinguir lesões potencialmente malignas (OPMD) de câncer já
estabelecido, orientando decisões terapêuticas distintas.

\subsection{Dataset}

Utilizamos o dataset SMART-OM~\cite{smartom}, composto por 2.469 imagens
intraorais coletadas de 331 pacientes em condições clínicas reais, com
smartphones Android e iOS, cobrindo oito regiões anatômicas da cavidade oral.
A Tabela~\ref{tab:dataset} apresenta a distribuição de imagens por classe
e a partição treino/teste adotada (80\%/20\%, sem estratificação por
paciente).

\begin{table}[ht]
  \centering
  \caption{Distribuição do dataset SMART-OM por classe.}
  \label{tab:dataset}
  \begin{tabular}{lrrrr}
    \toprule
    Classe & Total & Treino$^\dagger$ & Teste & \% do total \\
    \midrule
    Normal             & 2.145 & 1.716 & 429 & 86,9\% \\
    Variação do Normal &   179 &   143 &  36 &  7,2\% \\
    OPMD               &   125 &   100 &  25 &  5,1\% \\
    Câncer Oral        &    20 &    16 &   4 &  0,8\% \\
    \midrule
    \textbf{Total}     & \textbf{2.469} & \textbf{1.975} & \textbf{494} & 100\% \\
    \bottomrule
    \multicolumn{5}{l}{\footnotesize{$^\dagger$Antes de \textit{oversampling};
    classes minoritárias são superamostradas no treinamento (Seção~\ref{sec:metodologia}).}}
  \end{tabular}
\end{table}

O dataset é severamente desbalanceado: Normal representa 86,9\% das
imagens e Câncer Oral apenas 0,8\% (20 amostras), tornando a acurácia
global inadequada como critério principal de avaliação.


\subsection{Métricas de Avaliação}

Em conformidade com o benchmark SMART-OM~\cite{smartom}, reportamos para
cada tarefa e modelo as seguintes métricas sobre o conjunto de teste:

\begin{itemize}
  \item \textbf{Sensibilidade macro} (\textit{macro recall}): média não
    ponderada da sensibilidade por classe — métrica primária, pois penaliza
    igualmente erros em classes raras e frequentes;
  \item \textbf{F1-score macro}: média harmônica de precisão e sensibilidade,
    calculada por classe e depois agregada por média simples;
  \item \textbf{Acurácia global, Precisão e Especificidade} — macro e
    por classe;
  \item \textbf{AUC-ROC} macro e por classe, calculada via
    \textit{One-vs-Rest}.
\end{itemize}

A sensibilidade macro é adotada como critério de comparação principal
com o baseline, dado o desbalanceamento do dataset e a assimetria
clínica dos erros.

\section{Metodologia}
\label{sec:metodologia}

\subsection{Preparação do Dataset}

O dataset SMART-OM foi dividido aleatoriamente em 80\% para treinamento
e 20\% para teste, sem estratificação por paciente ou por região
intraoral, replicando o protocolo do benchmark de referência~\cite{smartom}.
Não foi utilizado conjunto de validação separado, dado o reduzido número
de amostras nas classes minoritárias — reservar uma terceira partição
reduziria ainda mais os dados disponíveis para treinamento das classes
OPMD (125 imagens) e Câncer Oral (20 imagens). E separar os dados do conjunto de treino tornaria as comparações inválidas para o protocolo de referência

Para mitigar o desbalanceamento severo (86,9\% Normal), aplicamos
\textit{oversampling} por duplicação aleatória nas classes minoritárias
do conjunto de treinamento, igualando a contagem de cada classe à maior
entre elas.
No conjunto de teste, nenhuma modificação é aplicada — as distribuições
originais são preservadas para avaliação.

\subsection{Arquiteturas Avaliadas}

Avaliamos a variante \textit{medium} de três famílias YOLO no modo de
classificação, todas inicializadas com pesos pré-treinados no
ImageNet e submetidas a \textit{fine-tuning} no SMART-OM:

\begin{itemize}
  \item \textbf{YOLOv8m-cls}~\cite{yolov8}: backbone CSPDarknet com
    módulos C2f;
  \item \textbf{YOLOv11m-cls}~\cite{yolo11}: evolução do YOLOv8 com
    módulos C3k2 e atenção espacial CBAM;
  \item \textbf{YOLO26m-cls}~\cite{yolo26}: arquitetura sem NMS,
    com \textit{ProgLoss} e otimizador MuSGD para convergência estável.
\end{itemize}

Todo o treinamento foi realizado via framework Ultralytics com
PyTorch~\cite{yolov8}, em GPU NVIDIA RTX 4060 Laptop (8 GB VRAM).

\subsection{Configuração de Treinamento}

A Tabela~\ref{tab:hiperparametros} resume os hiperparâmetros utilizados,
aplicados igualmente aos três modelos e às três tarefas de treinamento.

\begin{table}[ht]
  \centering
  \caption{Hiperparâmetros de treinamento.}
  \label{tab:hiperparametros}
  \begin{tabular}{ll}
    \toprule
    Parâmetro & Valor \\
    \midrule
    Épocas              & 100 \\
    Tamanho de batch    & 32 \\
    Tamanho de imagem   & 320 px \\
    Otimizador          & AdamW ($lr_0 = 10^{-3}$, $lr_f = 0{,}01$) \\
    \textit{Flip} horizontal & 0,5 \\
    Rotação             & $\pm$10° \\
    Translação          & 10\% \\
    Escala              & 30\% \\
    Jitter de matiz     & 0,005 \\
    Jitter de saturação & 0,7 \\
    Jitter de valor     & 0,4 \\
    \bottomrule
  \end{tabular}
\end{table}

\subsection{Tratamento do Desbalanceamento via Loss Ponderada}

Além do \textit{oversampling}, aplicamos ponderação de classes na função
de perda durante o treinamento. Os pesos são calculados pelo método de
frequência inversa normalizada:

\begin{equation}
  w_c = \frac{N}{C \cdot n_c}, \qquad
  w_c \leftarrow \frac{w_c}{\bar{w}}
  \label{eq:pesos}
\end{equation}

onde $N$ é o total de amostras de treinamento, $C$ é o número de
classes e $n_c$ é o número de amostras da classe $c$.
A normalização pela média $\bar{w}$ mantém a escala da loss equivalente
ao caso sem ponderação, preservando o \textit{learning rate} calibrado.

\subsection{Pipeline de Dois Estágios}

Como contribuição adicional, propomos um pipeline de classificação em
dois estágios sequenciais:

\textbf{Estágio 1 — Detector binário.}
Um classificador treinado na tarefa T\textsubscript{1} (Normal
vs.\ Anormal) recebe a imagem e decide se há alguma alteração da
mucosa. Caso a imagem seja classificada como Normal, o pipeline
encerra. Caso contrário, ela é encaminhada ao Estágio 2.

\textbf{Estágio 2 — Classificador de lesões.}
Um segundo classificador, treinado exclusivamente sobre as 760 imagens
anômalas (Variação do Normal, OPMD e Câncer Oral) após
\textit{oversampling}, distingue entre as três categorias patológicas.

A avaliação do pipeline é feita de forma \textit{end-to-end}: imagens
incorretamente classificadas como Normal pelo Estágio 1 nunca chegam ao
Estágio 2, sendo contabilizadas como erro na métrica final de quatro
classes. 
\section{Resultados e Discussão}
\label{sec:resultados}

\subsection{\texorpdfstring{Detecção Binária de Anomalias (T\textsubscript{1})}{Detecção Binária (T1)}}

A Tabela~\ref{tab:binary} apresenta os resultados completos na tarefa
binária, comparando o melhor modelo do benchmark SMART-OM
(ResNet18~\cite{smartom}) com os três modelos propostos.
Em negrito, o melhor valor por linha entre todos os modelos.

\begin{table}[htbp]
  \centering
  \caption{Desempenho na classificação binária
           (T\textsubscript{1}). Em negrito: melhor valor por linha.
           Baseline: ResNet18~\cite{smartom}.}
  \label{tab:binary}
  \small
  \begin{tabular}{llrrrr}
    \toprule
    \multicolumn{2}{l}{Métricas de Desempenho}
      & ResNet18 & YOLOv8m & YOLOv11m & YOLO26m \\
    \midrule
    \multicolumn{2}{l}{Acurácia (\%)}
      & 90,89 & 92,11 & 91,90 & \textbf{93,32} \\
    \midrule
    \multirow{5}{*}{Média Macro}
      & Precisão (\%)       & 79,74 & 84,20 & 83,22 & \textbf{87,31} \\
      & Sensibilidade (\%)  & 80,14 & 79,14 & 79,67 & \textbf{81,79} \\
      & Especificidade (\%) & 80,14 & 79,14 & 79,67 & \textbf{81,79} \\
      & F1-Score (\%)       & 79,94 & 81,37 & 81,30 & \textbf{84,24} \\
      & AUC                 & \textbf{0,91} & \textbf{0,91} & 0,88 & 0,87 \\
    \midrule
    \multirow{4}{*}{\parbox{1.6cm}{Métricas\\(Normal)}}
      & Precisão (\%)       & 94,87 & 94,32 & 94,51 & \textbf{95,00} \\
      & Sensibilidade (\%)  & 94,65 & 96,74 & 96,27 & \textbf{97,44} \\
      & Especificidade (\%) & 65,62 & 61,54 & 63,08 & \textbf{66,15} \\
      & F1-Score (\%)       & 94,76 & 95,51 & 95,38 & \textbf{96,20} \\
    \midrule
    \multirow{4}{*}{\parbox{1.6cm}{Métricas\\(Anormal)}}
      & Precisão (\%)       & 64,61 & 74,07 & 71,93 & \textbf{79,63} \\
      & Sensibilidade (\%)  & 65,62 & 61,54 & 63,08 & \textbf{66,15} \\
      & Especificidade (\%) & 94,65 & 96,74 & 96,27 & \textbf{97,44} \\
      & F1-Score (\%)       & 65,12 & 67,23 & 67,21 & \textbf{72,27} \\
    \bottomrule
  \end{tabular}
\end{table}

O YOLO26m-cls superou o ResNet18 em quase todas as métricas, com
destaque para precisão macro (+7,57 p.p.), F1-macro (+4,30 p.p.) e
sensibilidade da classe Anormal (66,15\% vs.\ 65,62\%).
O único indicador em que o baseline se mantém superior é a AUC (0,91
vs.\ 0,87), sugerindo que a discriminação global de ranking é
ligeiramente melhor no ResNet18, enquanto o YOLO26m-cls toma decisões
binárias mais precisas.
A Figura~\ref{fig:conf_binary} exibe a matriz de confusão do YOLO26m-cls:
97,44\% das imagens Normais são classificadas corretamente, com 33,85\%
de falsos negativos na classe Anormal.
A Figura~\ref{fig:sens_bars} compara a sensibilidade macro de todos os
modelos nas duas tarefas.

YOLOv8m-cls e YOLOv11m-cls mantiveram-se próximos ao ResNet18 em
sensibilidade macro (79,14\% e 79,67\%, respectivamente, contra 80,14\%
do baseline), com ganho de F1-macro em ambos os casos (+1,43 e
+1,36~p.p.).
O ganho de sensibilidade na classe Anormal obtido pelo YOLO26m-cls —
ainda que de apenas 0,53~p.p.\ em termos absolutos — tem relevância
clínica direta: em contexto de triagem, cada caso Anormal classificado
erroneamente como Normal representa um paciente que não é encaminhado
para avaliação especializada, podendo ter sua janela de intervenção
precoce comprometida.

\begin{figure}[htbp]
  \centering
  \includegraphics[width=0.52\textwidth]{fig_confusion_binary.pdf}
  \caption{Matriz de confusão do YOLO26m-cls na tarefa binária
           (T\textsubscript{1}).}
  \label{fig:conf_binary}
\end{figure}

\begin{figure}[htbp]
  \centering
  \includegraphics[width=\textwidth]{fig_sensitivity.pdf}
  \caption{Sensibilidade macro dos modelos propostos vs.\ baseline do
           SMART-OM nas tarefas binária (a) e multiclasse (b).
           A linha tracejada indica o melhor resultado do benchmark.}
  \label{fig:sens_bars}
\end{figure}

\FloatBarrier
\subsection{\texorpdfstring{Classificação Multiclasse (T\textsubscript{2})}{Classificação Multiclasse (T2)}}

A Tabela~\ref{tab:multiclass} detalha os resultados completos por classe
na tarefa de quatro categorias, tendo como baseline o ResNet50~\cite{smartom}
— o melhor modelo reportado pelo benchmark original nessa tarefa.
A classificação multiclasse é substancialmente mais desafiadora que a
tarefa binária: além do desbalanceamento geral, as classes OPMD (5,1\%)
e Câncer Oral (0,8\%) apresentam características visuais sutis e contagens
de teste reduzidas, tornando a sensibilidade macro o critério mais
informativo para comparação entre modelos.

\begin{table}[tp]
  \centering
  \caption{Desempenho na classificação multiclasse
           (T\textsubscript{2}). Em negrito: melhor valor por linha.
           Baseline: ResNet50~\cite{smartom}.}
  \label{tab:multiclass}
  \small
  \begin{tabular}{llrrrr}
    \toprule
    \multicolumn{2}{l}{Métricas de Desempenho}
      & ResNet50 & YOLOv8m & YOLOv11m & YOLO26m \\
    \midrule
    \multicolumn{2}{l}{Acurácia (\%)}
      & 89,07 & \textbf{91,30} & 90,08 & 90,69 \\
    \midrule
    \multirow{5}{*}{Média Macro}
      & Precisão (\%)       & 74,62 & \textbf{81,41} & 77,86 & 77,13 \\
      & Sensibilidade (\%)  & 61,57 & 57,06 & 52,94 & \textbf{69,57} \\
      & Especificidade (\%) & 84,63 & 89,42 & 87,11 & \textbf{89,92} \\
      & F1-Score (\%)       & 65,99 & 61,36 & 58,80 & \textbf{71,46} \\
      & AUC                 & 0,91  & 0,92  & 0,92  & \textbf{0,93}  \\
    \midrule
    \multirow{4}{*}{\parbox{1.6cm}{Métricas\\(Normal)}}
      & Precisão (\%)       & 91,83 & 94,38 & 93,10 & \textbf{94,75} \\
      & Sensibilidade (\%)  & 96,74 & \textbf{97,90} & 97,44 & 96,74 \\
      & Especificidade (\%) & 42,19 & 61,54 & 52,31 & \textbf{64,62} \\
      & F1-Score (\%)       & 94,22 & \textbf{96,11} & 95,22 & 95,73 \\
    \midrule
    \multirow{4}{*}{\parbox{1.6cm}{Métricas\\(Variação)}}
      & Precisão (\%)       & 40,00 & \textbf{75,00} & 60,00 & 57,89 \\
      & Sensibilidade (\%)  & 16,22 & \textbf{33,33} & \textbf{33,33} & 30,56 \\
      & Especificidade (\%) & 98,03 & \textbf{99,13} & 98,25 & 98,25 \\
      & F1-Score (\%)       & 23,08 & \textbf{46,15} & 42,86 & 40,00 \\
    \midrule
    \multirow{4}{*}{\parbox{1.6cm}{Métricas\\(OPMD)}}
      & Precisão (\%)       & \textbf{66,67} & 56,25 & 58,33 & 55,88 \\
      & Sensibilidade (\%)  & 66,67 & 72,00 & 56,00 & \textbf{76,00} \\
      & Especificidade (\%) & \textbf{98,30} & 97,01 & 97,87 & 96,80 \\
      & F1-Score (\%)       & \textbf{66,67} & 63,16 & 57,14 & 64,41 \\
    \midrule
    \multirow{4}{*}{\parbox{1.6cm}{Métricas\\(Câncer Oral)}}
      & Precisão (\%)       & \textbf{100,00} & \textbf{100,00} & \textbf{100,00} & \textbf{100,00} \\
      & Sensibilidade (\%)  & 66,67 & 25,00 & 25,00 & \textbf{75,00} \\
      & Especificidade (\%) & \textbf{100,00} & \textbf{100,00} & \textbf{100,00} & \textbf{100,00} \\
      & F1-Score (\%)       & 80,00 & 40,00 & 40,00 & \textbf{85,71} \\
    \bottomrule
  \end{tabular}
\end{table}

O YOLO26m-cls foi o único modelo proposto a superar o ResNet50 na
métrica primária, com sensibilidade macro de 69,57\% contra 61,57\%
do baseline (+8,0~p.p.), além de F1-macro superior (+5,47~p.p.) e
maior AUC (0,93 vs.\ 0,91).
YOLOv8m-cls e YOLOv11m-cls ficaram abaixo do ResNet50 em sensibilidade
macro (57,06\% e 52,94\%, respectivamente), demonstrando que o ganho
não é uniforme na família YOLO e está concentrado na arquitetura mais
recente.
Na classe OPMD — clinicamente crítica por representar a janela de
intervenção pré-maligna — o YOLO26m-cls atingiu 76,00\% de
sensibilidade contra 66,67\% do baseline.

\begin{figure}[htbp]
  \centering
  \includegraphics[width=0.7\textwidth]{fig_confusion_multiclass.pdf}
  \caption{Matriz de confusão do YOLO26m-cls na tarefa multiclasse
           (T\textsubscript{2}).}
  \label{fig:conf_multi}
\end{figure}

\FloatBarrier
\subsection{\texorpdfstring{Pipeline de Dois Estágios (T\textsubscript{2})}{Pipeline de Dois Estágios (T2)}}

A Tabela~\ref{tab:twostage} compara o pipeline de dois estágios com as
abordagens \textit{one-stage} e o baseline do SMART-OM na tarefa de
classificação multiclasse.
O pipeline combina o melhor detector binário (YOLO26m-cls) com o melhor
classificador de lesões (YOLOv8m-cls no estágio 2), selecionados
individualmente por sensibilidade macro em seus respectivos conjuntos de
validação.

\begin{table}[htbp]
  \centering
  \caption{Comparação entre abordagens na tarefa multiclasse (T\textsubscript{2}).
           Em negrito: melhor valor por linha.
           Two-Stage: YOLO26m-cls (binário) + YOLOv8m-cls (lesões).}
  \label{tab:twostage}
  \footnotesize
  \begin{tabular}{llrrrrr}
    \toprule
    \multicolumn{2}{l}{Métricas de Desempenho}
      & ResNet50 & YOLOv8m & YOLOv11m & YOLO26m & Two-Stage \\
    \midrule
    \multicolumn{2}{l}{Acurácia (\%)}
      & 89,07 & 91,30 & 90,08 & 90,69 & \textbf{91,90} \\
    \midrule
    \multirow{3}{*}{Média Macro}
      & Sensibilidade (\%)  & 61,57 & 57,06 & 52,94 & 69,57 & \textbf{71,83} \\
      & F1-Score (\%)       & 65,99 & 61,36 & 58,80 & 71,46 & \textbf{74,94} \\
      & AUC                 & 0,91  & 0,92  & \textbf{0,92}  & \textbf{0,93}  & 0,83  \\
    \midrule
    \multirow{4}{*}{\parbox{1.8cm}{Sensibilidade\\por Classe (\%)}}
      & Normal      & 96,74 & \textbf{97,90} & 97,44 & 96,74 & \textbf{97,44} \\
      & Variação    & 16,22 & \textbf{33,33} & \textbf{33,33} & 30,56 & 38,89 \\
      & OPMD        & 66,67 & 72,00 & 56,00 & \textbf{76,00} & \textbf{76,00} \\
      & Câncer Oral & 66,67 & 25,00 & 25,00 & \textbf{75,00} & \textbf{75,00} \\
    \bottomrule
    \multicolumn{7}{l}{\footnotesize{Baseline ResNet50 de~\cite{smartom}; modelos YOLO treinados neste trabalho.}}
  \end{tabular}
\end{table}

O pipeline de dois estágios obteve a maior sensibilidade macro (71,83\%)
e o maior F1-macro (74,94\%) entre todas as abordagens avaliadas,
superando o melhor modelo \textit{one-stage} (YOLO26m, 69,57\%) em
2,26~p.p.\ e o baseline em 10,26~p.p.
A Variação do Normal registrou a maior sensibilidade de todos os métodos
(38,89\%), beneficiada pela especialização do Estágio~2 em distinguir
apenas entre categorias patológicas.

A AUC do pipeline (0,83) ficou substancialmente abaixo dos modelos
\textit{one-stage} (0,92--0,93).
Essa queda é estrutural: a AUC avalia a qualidade do \textit{ranking}
de probabilidades, e no pipeline de dois estágios as probabilidades dos
dois estágios não são calibradas em escala comum.
% Essa descontinuidade entre as duas escalas de confiança degrada a curva
% ROC, sem que haja perda equivalente na qualidade das decisões binárias
% tomadas em cada estágio.

\FloatBarrier
\subsection{Discussão}

Os resultados revelam um panorama nuançado. O YOLO26m-cls é o modelo com
melhor desempenho \textit{one-stage} nas métricas clinicamente mais
relevantes — sensibilidade macro (+8,0 p.p.\ sobre o ResNet50) e
F1-macro (+5,47 p.p.).
O pipeline de dois estágios vai além, atingindo 71,83\% de sensibilidade
macro e 74,94\% de F1-macro — os melhores valores entre todas as
abordagens avaliadas — ao custo de maior complexidade de implantação e
de uma AUC degradada pela falta de calibração entre estágios.
YOLOv8m-cls e YOLOv11m-cls ficam abaixo do ResNet50 em sensibilidade
macro \textit{one-stage}, evidenciando que o ganho não é universal à
família YOLO, mas concentrado na arquitetura YOLO26m.

Na tarefa binária, o YOLO26m-cls também lidera, superando o ResNet18
em todas as métricas exceto AUC (0,87 vs.\ 0,91). Isso indica que o
modelo produz probabilidades menos calibradas para ranking, mas toma
decisões de limiar com maior precisão.

A classe \textit{Variação do Normal} permanece o principal gargalo em
todos os modelos avaliados, incluindo o baseline (sensibilidade de
16--43\% no \textit{one-stage}). Suas imagens partilham características
visuais com as demais classes, tornando-a intrinsecamente difícil de
separar sem informação adicional. A classe \textit{Câncer Oral}, com
apenas 4 amostras de teste, apresenta alta variância — o YOLO26m-cls e
o pipeline acertam 3/4 (75\%) esses números devem ser interpretados com
cautela dado o tamanho ínfimo da partição.


\section{Conclusão}
\label{sec:conclusao}

Este trabalho avaliou três arquiteturas YOLO modernas no modo de
classificação — YOLOv8m-cls, YOLOv11m-cls e YOLO26m-cls — sobre o
dataset público SMART-OM~\cite{smartom}, em duas tarefas: detecção
binária de anomalias (T\textsubscript{1}) e classificação multiclasse
de quatro categorias diagnósticas (T\textsubscript{2}).
O critério de sucesso estabelecido nos objetivos foi atingido: o
YOLO26m-cls superou o benchmark original em ambas as tarefas,
alcançando 66,15\% de sensibilidade na classe Anormal contra 65,62\%
do ResNet18 na tarefa binária, e 69,57\% de sensibilidade macro contra
61,57\% do ResNet50 na tarefa multiclasse — um ganho de 8,0 p.p.\ na
métrica primária.
O F1-macro também melhorou em ambas as tarefas (+4,30 p.p.\ em
T\textsubscript{1} e +5,47 p.p.\ em T\textsubscript{2}).

Esses resultados demonstram que arquiteturas concebidas para detecção
em tempo real, quando adaptadas à tarefa de classificação e treinadas
com ponderação de classes, são competitivas com CNNs clássicas em
cenários clínicos desbalanceados.
A ponderação por frequência inversa e o \textit{oversampling} mostraram-se
particularmente eficazes para a classe Câncer Oral
(sensibilidade de 75\% com YOLO26m-cls contra 66,67\% do ResNet50),
apesar das apenas 4 amostras de teste disponíveis.
Contudo, a família YOLO não apresentou ganhos uniformes: YOLOv8m-cls e
YOLOv11m-cls ficaram abaixo do ResNet50 em sensibilidade macro na
tarefa multiclasse, e o YOLO26m-cls perdeu em AUC na tarefa binária
(0,87 vs.\ 0,91), indicando que a calibração de probabilidades
permanece um ponto de atenção.

O trabalho apresenta limitações que devem ser consideradas na
interpretação dos resultados.
A classe Câncer Oral conta com apenas 20 imagens no total e 4 no
conjunto de teste, tornando suas métricas altamente sensíveis a
predições individuais.
A classe Variação do Normal permanece o principal gargalo — com
sensibilidade entre 16\% e 43\% em todos os modelos avaliados, inclusive
no baseline — sugerindo que suas características visuais se sobrepõem
às demais classes de forma intrínseca.

Como trabalhos futuros, destacam-se: (i) adotar validação cruzada
estratificada para produzir estimativas de desempenho mais robustas
e comparáveis ao benchmark; (ii) investigar estratégias de
\textit{data augmentation} específicas para imagens intraorais, como
transformações de cor controladas que preservem as propriedades
diagnósticas da mucosa; e (iii) explorar técnicas de explicabilidade
— como mapas de ativação Grad-CAM — para identificar as regiões da
imagem que sustentam as predições e aumentar a confiança clínica
nos modelos.

\bibliographystyle{sbc}
\bibliography{referencias}

\end{document}



